\itemhead{习题2.1 · 第4题（教材p.43）}{XJ-B1-C02-HW-P043-Y0819}
\label{XJ-B1-C02-HW-P043-Y0819}
{\footnotesize 湘教版教材p.43 / PDF第50页 / 习题2.1 / 第4题}
\field{原题}{下列结论是否成立？若成立，试说明理由；若不成立，试举出反例。
\subq{1}如果 $c-a>c-b$，那么 $a<b$；
\subq{2}若 $ab>c,b>0$，则 $a>\dfrac cb$；
\subq{3}若 $ac>bc$，则 $a>b$；
\subq{4}若 $a>b,c>d$，则 $a-c>b-d$。}
\field{参考答案}{（1）（2）成立；（3）不成立，取 $a=1,b=2,c=-1$；（4）不成立，取 $a=2,b=1,c=4,d=1$。}
{\small \field{关联知识}{I02 不等式基本性质与符号条件（教材帮p.45、46 / PDF54、55）}
\field{方法与易错点}{每一步明确所乘除因子的正负与非零，或作差求证。 易错：未知符号就直接乘除；将 $c^2$ 当作必为正数而忘记 $c=0$。}
\field{答案依据}{答案册相关题组：PDF第12、13页（册内页11、12）；按当前题库原题逐项重算，不按旧版题号机械对套。}}
